\section{Experimental Setup: Full Details}
\label{app:setup}

This appendix expands the condensed Setup in \S\ref{sec:experiments:setup} with the formal metric definitions and the runtime infrastructure.

\paragraph{Seed agent.}
The seed configuration, denoted \textbf{NexAU\textsubscript{0}}, is a simple code agent built on the NexAU framework of \S\ref{sec:method:harness} that exposes only the \texttt{bash} tool to the model, with no skills, no middleware, and no long-term memory. Every iteration of the AHE outer loop edits this workspace, so all reported gains are measured against NexAU\textsubscript{0} as the common starting point.

\paragraph{Runtime infrastructure.}
All runs use the NexAU framework of \S\ref{sec:method:harness} to instantiate the coding agent. Harbor dispatches tasks, isolates each rollout, and verifies pass/fail, with a 3600-second per-task timeout. Every rollout runs inside a fresh E2B remote sandbox, so shell side-effects cannot leak between tasks. \texttt{InMemoryTracer} records trajectories and mirrors them to Langfuse. The Agent Debugger executes at concurrency 16 with a 600-second per-task timeout.

\paragraph{Hyperparameters.}
Table~\ref{tab:hparams} consolidates the operating points of the four agents and the outer loop, taken from the configuration snapshot of the reference run.

\begin{table}[h]
    \centering
    \small
    \setlength{\tabcolsep}{6pt}
    \renewcommand{\arraystretch}{1.05}
    \caption{Hyperparameters of the reference AHE run on Terminal-Bench~2, taken from the configuration snapshot of the run that produced the AHE numbers in Section~\ref{sec:experiments:main}. All four agents share GPT-5.4 as the backing model and differ only in reasoning tier, sampling, and per-component limits.}
    \label{tab:hparams}
    \begin{tabular}{lll}
        \toprule
        Block & Hyperparameter & Value \\
        \midrule
        Code Agent
            & Backing model              & GPT-5.4 \\
            & Reasoning effort           & high \\
            & Temperature                & 0.7 \\
            & Top-$p$                    & 0.95 \\
            & Max context                & 200{,}000 tokens \\
            & Max generation per turn    & 32{,}000 tokens \\
            & Max model turns            & 300 \\
        \midrule
        Agent Debugger
            & Backing model              & GPT-5.4 \\
            & Concurrent tasks           & 16 \\
            & Per-task timeout           & 600\,s \\
            & Retry attempts             & 3 \\
            & Retry backoff              & 2\,s \\
        \midrule
        Evolve Agent
            & Backing model              & GPT-5.4 \\
            & Reasoning effort           & xhigh \\
            & Temperature                & 0.7 \\
            & Max context                & 200{,}000 tokens \\
            & Max generation per turn    & 32{,}000 tokens \\
            & Max model turns            & 500 \\
            & Skill packages             & 3 \\
            & Context compaction threshold & 0.75 \\
        \midrule
        Explore Agent
            & Backing model              & GPT-5.4 \\
            & Reasoning effort           & xhigh \\
            & Wall-clock budget          & 60\,min \\
            & Web sources                & 10 \\
            & Code sources               & 1 \\
        \midrule
        Outer loop / dispatcher
            & Iterations $T$             & 10 \\
            & Rollouts per task $k$      & 2 \\
            & Concurrent rollouts        & 96 \\
            & Sandbox                    & E2B remote \\
            & Sandbox lifetime           & 3{,}600\,s \\
            & Dataset                    & Terminal-Bench~2, 89 tasks \\
        \bottomrule
    \end{tabular}
\end{table}


\paragraph{Terminal-bench difficulty labels.}
The official terminal-bench-2 leaderboard\footnote{\url{https://www.tbench.ai/benchmarks/terminal-bench-2}} partitions the 89-task subset into 4 easy, 55 medium, and 30 hard tasks.

\paragraph{pass@1.}
For a configuration on a task set $D$ with $k$ rollouts per task, let $r_{i,j} \in \{0, 1\}$ denote the binary reward of rollout $j$ on task $i$. The pass@1 score is the mean
\begin{equation}
    \mathrm{pass@1} = \frac{1}{k|D|} \sum_{i=1}^{|D|} \sum_{j=1}^{k} r_{i,j}.
\end{equation}
Trials that terminate on an infrastructure exception, such as a sandbox crash or API timeout, contribute $r=0$ rather than being dropped, a strictly harsher convention than discarding failures that keeps our numbers comparable to the official terminal-bench leaderboard. The rollout count $k$ varies across experiments; each table states it explicitly.

\paragraph{Token cost and Succ/Mtok.}
For token cost we count every LLM call as prompt plus completion across the rollout and report the mean over completed trials in thousands, denoted Tokens k; infrastructure-aborted trials are excluded to avoid truncated figures. To compare configurations that trade accuracy for cost we combine the two via
\begin{equation}
    \mathrm{Succ/Mtok} = \frac{\mathrm{pass@1} \times 10^{6}}{\mathrm{mean\ tokens\ per\ trial}},
\end{equation}
the expected number of successes per million tokens. The main paper reports pass@1 and Tokens k separately so each axis stays legible; Table~\ref{tab:swe-verified-succ-mtok} folds them into Succ/Mtok per repository on SWE-bench-verified, derived from the pass@1 and Tokens k columns of Table~\ref{tab:swe-verified-analysis}.

\begin{table}[t]
    \centering
    \small
    \setlength{\tabcolsep}{6pt}
    \renewcommand{\arraystretch}{1.05}
    \caption{Cost-efficiency on SWE-bench-verified, reported as Succ/Mtok, the expected successes per million tokens. Values are derived from Table~\ref{tab:swe-verified-analysis} as $\mathrm{pass@1} \times 10^{3} / \text{Tokens k}$. Higher is better. Per-row bold marks the best.}
    \label{tab:swe-verified-succ-mtok}
    \begin{tabular}{lccccc}
        \toprule
        Repo & $N$ & ACE & TF-GRPO & NexAU\textsubscript{0} & \textbf{AHE} \\
        \midrule
        All            & 500 & 1.10 & 1.27 & 1.43 & \textbf{1.64} \\
        \midrule
        django         & 231 & 1.12 & 1.35 & 1.50 & \textbf{1.67} \\
        sympy          &  75 & 1.15 & 1.19 & 1.43 & \textbf{1.48} \\
        sphinx-doc     &  44 & 0.62 & 0.78 & 0.93 & \textbf{1.07} \\
        matplotlib     &  34 & 1.14 & 1.33 & 1.51 & \textbf{1.88} \\
        scikit-learn   &  32 & 2.08 & 2.48 & 3.06 & \textbf{3.40} \\
        pydata         &  22 & 1.37 & 1.50 & 2.00 & \textbf{2.15} \\
        astropy        &  22 & 1.08 & 1.26 & 0.82 & \textbf{1.81} \\
        \bottomrule
    \end{tabular}
\end{table}

